\begin{table}[!htbp]
\centering
\small
\setlength{\tabcolsep}{4pt}
\renewcommand{\arraystretch}{1.14}
\caption{Raw-request lifecycle states at specified checkpoints (S8). Each row was observed in all six runs of that scenario: three repetitions of each information-matched implementation. Counts were reconstructed from the saved client and sink databases. These are checkpoint observations, not continuous measurements or independent fault-rate samples.}
\label{tab:lifecycle-states}
\begin{tabularx}{\linewidth}{@{}>{\raggedright\arraybackslash}Xrrrr@{}}
\toprule
Scenario checkpoint & \shortstack{Registered\\permits} & \shortstack{Held\\units} & \shortstack{Sink\\effects} & \shortstack{Local\\receipts}\\
\midrule
After distinct-request contention & 8 & 0 & 2 & 2 \\
After unsupported-claim release & 1 & 0 & 0 & 0 \\
After interrupted issuance rollback & 0 & 0 & 0 & 0 \\
Early cleanup before arm & 1 & 1 & 0 & 0 \\
After expired unarmed cancellation & 1 & 0 & 0 & 0 \\
Cleanup after committed effect and timeout & 1 & 1 & 1 & 0 \\
\bottomrule
\end{tabularx}
\par\smallskip
\begin{minipage}{\linewidth}\small
Registered permits include retained non-action and cancelled records. The last row is allowed: a missing receipt does not justify releasing the held unit after the sink effect has committed. All 36 formal runs met the declared checks. Each run finished with two effects and two receipts, except the mixed-claim scenario, which finished with three of each (two alerts and one review), giving 78 effects and 78 receipts overall. Each final state had zero held units. The audit retained 36 direct final database pairs and 126 checkpoint backup pairs: 90 non-final snapshots and 36 final-state backups.
\end{minipage}
\end{table}
